%% ===================================================================
\subsection{Template design}%
\label{sec:template-design}
%% ===================================================================

\FractalShark{} makes pervasive use of C++ templates to select numeric
types, iteration widths, and algorithmic variants at compile time rather
than at run time.  Two considerations motivate this choice.  First,
the inner loops of both reference-orbit computation and per-pixel
perturbation execute billions of arithmetic operations; virtual dispatch
or run-time branch mispredictions on every iteration would impose
measurable overhead.  Second, encoding precision contracts in the type
system allows the compiler to verify that high-precision reference data
is never silently truncated when it is handed to a lower-precision
perturbation kernel, catching a class of errors that would otherwise
manifest only as subtle rendering artifacts.

The template parameters fall into four largely orthogonal dimensions,
described in the following subsections.


\subsubsection{Iteration-counter width (\code{IterType})}
\label{subsec:template-itertype}

Every orbit and perturbation result is indexed by an iteration counter
whose type is either \code{uint32\_t} or \code{uint64\_t}.  A 32-bit
counter suffices for the vast majority of renders, supporting up to
$2^{32}-1 \approx 4.3\times10^{9}$ iterations.  Certain ultra-deep
locations, however, require iteration limits that exceed this range.
Switching to a 64-bit counter doubles the per-element storage cost for
every orbit and iteration buffer, so \FractalShark{} treats the counter
width as a compile-time choice rather than unconditionally paying the
wider cost.  The \code{IterType} parameter threads through
\code{PerturbationResults}, \code{LAReference}, and the rendering
kernels, ensuring that all components agree on the iteration index
representation.


\subsubsection{Reference-orbit type (\code{T})}
\label{subsec:template-t}

The parameter~\code{T} selects the numeric type used to store the
reference orbit $\{z_n^\star\}$ and, by extension, the type in which the
high-precision reference computation is performed.  It determines both
the precision available during orbit iteration and the storage
representation uploaded to the GPU for use by perturbation kernels.
The principal choices are:

\begin{itemize}
  \item \code{float} and \code{double}---IEEE~754 single and double
        precision, used at shallow zoom depths where the full mantissa
        suffices to represent the orbit accurately.
  \item \code{CudaDblflt<MattDblflt>}---a double-single representation
        that pairs two 32-bit floats to obtain roughly double the
        mantissa bits of \code{float} alone, with natively vectorisable
        GPU arithmetic (\cref{sec:extended-precision-types}).
  \item \code{HDRFloat<double>} and \code{HDRFloat<float>}---a
        high-dynamic-range wrapper that adds an explicit exponent to the
        mantissa, preventing underflow in the deep-zoom regime where
        orbit values span hundreds of orders of magnitude
        (\cref{sec:hdr-float}).
  \item \code{HDRFloat<CudaDblflt<MattDblflt>>}---the composition of
        the HDR exponent with the double-single mantissa, providing
        both extended range and extended precision.
  \item \code{HpSharkFloat<SharkFloatParams>}---the GPU-resident
        arbitrary-precision type used for reference-orbit computation
        at extreme zoom depths (\cref{subsec:ref-orbit-gpu}).
\end{itemize}

\noindent
The type~\code{T} appears as the second template parameter of
\code{PerturbationResults<IterType,\allowbreak T,\allowbreak PExtras>}
and as the orbit-storage type inside \code{RefOrbitCalc}.


\subsubsection{Perturbation subtype (\code{SubType})}
\label{subsec:template-subtype}

While~\code{T} governs the precision of the shared reference orbit,
per-pixel perturbation operates at a lower precision because each
pixel evolves only a small delta $\Delta z_n$ relative to $z_n^\star$.
The \code{SubType} parameter selects this lower-precision type.  It
is typically \code{float}, \code{double}, or
\code{CudaDblflt<MattDblflt>}, and for HDR-wrapped references the
corresponding \code{HDRFloat} variant.  The mapping from \code{T} to
\code{SubType} is encoded by the \code{SubTypeChooser} trait: when
\code{T} is a plain scalar the sub-type is the same scalar; when
\code{T} is an \code{HDRFloat<S>} the sub-type extracts the
underlying~\code{S}.

\code{SubType} appears as an explicit template parameter on
\code{LAReference<IterType,\allowbreak T,\allowbreak SubType,\allowbreak
PExtras>} and is threaded into the perturbation kernels, where it
controls the precision of the per-pixel inner loop.


\subsubsection{Perturbation extras (\code{PExtras})}
\label{subsec:template-pextras}

The \code{Perturb\-Extras} enumeration is a compile-time flag attached to
every \code{Perturbation\-Results} instantiation.  It controls whether
the orbit stores auxiliary data beyond the orbit values themselves:

\begin{itemize}
  \item \code{Disable}---store only the orbit; no extra bookkeeping.
  \item \code{Bad}---record \emph{glitch} information (iterations where
        the perturbation approximation loses accuracy) alongside the
        orbit, enabling downstream glitch-correction passes.
  \item \code{SimpleCompression} and \code{MaxCompression}---store the
        metadata required for reference-orbit compression
        (\cref{sec:ref-compression-zhuoran}).
\end{itemize}

\noindent
Making this a template parameter rather than a run-time flag eliminates
the storage overhead of the auxiliary arrays in the common case where
they are not needed, and allows the compiler to elide the associated
bookkeeping code entirely.


\subsubsection{Instantiation combinatorics}
\label{subsec:template-combinatorics}

The Cartesian product of the four dimensions yields a large number of
potential type combinations.  With two \code{IterType} widths, six
representative \code{T} values, and three commonly used
\code{PExtras} settings (\code{Disable}, \code{Bad},
\code{Simple\-Compression}), the \code{PerturbationResults} class
alone accounts for $2\times 6\times 3 = 36$ concrete instantiations.
This count is visible in the \code{AwesomeVariant} type alias defined
in \code{RefOrbitCalc.h}, which enumerates all 36~alternatives as a
\code{std::variant} so that perturbation results of any type
combination can be stored and dispatched uniformly at run time.

Additional template parameters (\code{SubType}, benchmark modes, reuse
modes) on related classes further multiply the instantiation count.  The
full set of explicit instantiations is available in
\code{RefOrbitCalcTemplates.h}.


\subsubsection{Explicit instantiation strategy}
\label{subsec:template-explicit-inst}

Because every translation unit that includes a template header would
otherwise independently instantiate each specialization it references,
\FractalShark{} relies on \emph{explicit template instantiation} to
centralize code generation and control compile times.

For the perturbation pipeline, the header
\code{RefOrbitCalcTemplates.h} defines a family of preprocessor macros
(\code{InstantiateIsPerturbation\-Result\-UsefulHere},
\code{InstantiateAddPerturbation\-Reference\-Point}, and others) that
expand to explicit instantiation declarations for every supported
\code{(IterType, T, SubType, PExtras)} combination.  Each macro accepts
the type arguments, expands to the required \code{template} declaration,
and is invoked once per combination immediately below its definition.
This pattern confines the instantiation cost to a single translation
unit while allowing other files to use only the declarations.

For the GPU arithmetic subsystem, a separate code-generation tool
(\code{HpShark\-Instantiate}) produces \code{.cu} files containing
explicit instantiations of \code{HpSharkFloat}, the NTT multiply
kernels, the reference-orbit GPU loop, and associated helper templates.
The generator partitions the output into paired
files---one for production parameter sets
(\code{SharkParams1}\,--\,\code{12}),
one for non-periodicity variants
(\code{SharkParamsNP1}\,--\,\code{NP12})---and groups
instantiations into named \emph{batches} (e.g.\
\code{HpSharkFloat\_Conversions}, \code{MultiplyNTT},
\code{HpSharkReference}) so that each batch compiles as an
independent translation unit.  This decomposition enables parallel
compilation and limits the peak memory consumption of the CUDA
compiler.


\subsubsection{The \code{SharkFloatParams} pattern}
\label{subsec:template-sharkparams}

The GPU high-precision arithmetic library takes the compile-time
parameterization one step further by encoding not only type choices but
also \emph{algorithmic configuration} in a template parameter.
The struct
\code{GenericSharkFloatParams<pNumDigits,\allowbreak Periodicity,\allowbreak
NewtonRaphson,\allowbreak SubTypeT>}
bundles several \code{static constexpr} members that the CUDA compiler
can fold into kernel code:

\begin{itemize}
  \item \code{GlobalNumUint32}---the number of 32-bit limbs in each
        high-precision operand, derived from \code{pNumDigits}.  Because
        the compiler knows this value at template instantiation time,
        loop trip counts in the NTT multiply and carry-propagation
        kernels are compile-time constants, enabling aggressive
        unrolling and register allocation.
  \item \code{EnablePeriodicity}---controls whether the reference-orbit
        kernel checks for period detection at each iteration.  Disabling
        it in the non-periodicity (\code{SharkParamsNP*}) and
        Newton--Raphson (\code{SharkParamsNR*}) variants eliminates
        the associated branch and storage overhead.
  \item \code{EnableNewtonRaphson}---when set, the kernel additionally
        tracks the derivative $d_n = \partial z_n/\partial c$ alongside
        the orbit, allocating extra GPU scratch memory and executing the
        derivative recurrence at each iteration.  This derivative is required by
        the Feature Finder's Newton's method solver
        (\cref{sec:feature-finder}).
  \item \code{NTTPlan}---a compile-time NTT plan computed by
        \code{SharkNTT::BuildPlanPrime} from \code{GlobalNumUint32},
        fixing the NTT transform length and modular-arithmetic
        parameters at compile time.
\end{itemize}

\noindent
Concrete aliases such as
\code{SharkParams1}\,--\,\code{SharkParams12} fix
\code{pNumDigits} to powers of two from $256$ through $524{,}288$
(providing roughly $2{,}500$ to $5{,}000{,}000$ decimal digits of
precision) with periodicity enabled and Newton--Raphson disabled.
Additional alias families---\code{SharkParamsNP*} (no periodicity),
\code{SharkParamsNR*} (Newton--Raphson enabled),
\code{SharkParamsDbl*} (\code{double} sub-type), and
\code{SharkParamsDbf*} (\code{CudaDblflt} sub-type)---provide the
remaining feature-flag and sub-type permutations.  Each alias is
independently instantiated by the code generator described in
\cref{subsec:template-explicit-inst}, yielding $12\times 5 = 60$
parameter-set instantiations per GPU kernel template.

The remaining sections collect operational engineering references for tests,
kernels, display output, orbit interchange, and portability.
